%*******************************************************
% Preface
%*******************************************************
\pdfbookmark[1]{Preface}{preface}
\chapter*{Preface}
\markboth{\spacedlowsmallcaps{Preface}}{\spacedlowsmallcaps{Preface}}

Every particle carries two numbers. One tells how much matter it holds: its mass. The other tells how strongly it takes part in electricity: its charge. Physics has always written these two numbers apart, each in its own law, each with its own constant. This book asks a simple question with a long history behind it: could one number hold them both?

The question did not arrive all at once. It grew over many years, out of three pieces of work that only make full sense when they are read together.

The first piece is the \emph{General Sinus}. It begins with a small doubt about something everyone learns at school. The sine and the cosine are defined on a right triangle. But why a right triangle? Most triangles in the world are not right-angled. The General Sinus gives the sine and the cosine a second variable, so that they can measure any triangle, and it gives back the familiar functions when the angle happens to be right \citep{Derraz2026GSP}.


The second piece is the \emph{General Numbers}. The General Sinus does more than add a variable to two old functions. It frees trigonometry from a single point of view, the vertical projection, and lets us look at the same point from any projection angle we choose. That freedom of perspective is what unlocks the complex plane. For about two centuries the complex plane has been built on one fixed frame, with its second axis standing at a right angle to the first. Once the projection is free, the frame is free too: the second axis may lean at any angle, and its unit may take a length of its own. This is the General Number Plane. It extends the complex plane and its frame, and it keeps the complex numbers inside it, as the one frame in which the second axis stands upright. The mathematics of the General Numbers is also published on its own, as a paper with complete proofs \citep{DerrazGNPreprint2026}.

The third piece is the \emph{Cosmic Mass}. Inside the General Number Plane there is room for a number whose first part is a mass and whose second part is a charge. That number is the cosmic mass. With it, a single law of interaction can be written for two particles, and that law contains the gravity of Newton, a term shaped like the electrostatics of Coulomb, and a third term that neither of them contains. Read by its real part, the law gives back the laws of Newton and Coulomb exactly, and the third term adds no force between the two particles.

The heart of the book is Chapter~4, on General Numbers, and Chapter~6, on the Interaction Law. The chapters before them are there to carry you to them and to make them easy to read: a first chapter that shows you the destination at once, a second on how a new theory is judged and on what this volume shows, a chapter on the General Sinus, and a prologue on why the language of physics must sometimes grow. The chapters after them put the cosmic mass to work. Chapter~7 draws every particle, ion and nucleus in one frame; Chapter~8 reads the mass of a nucleus off its neighbours, and the energy of a fusion before anyone weighs it; and Chapter~9 writes, for slow bodies, one law of fields and motion that holds Newton's gravity, Maxwell's equations and the force of Lorentz. Chapter~10, in Part~\ref{pt:open}, opens the question that the second volume takes up: what the third term of the law could mean, and what measurement already says about it. Chapter~11 draws the whole frame: antimatter, the half of negative mass, and the black holes. Short scenes from the long history of numbers, angles, matter and force run through them all.

This is the first of two volumes. It shows that the language of the cosmic mass can speak the old laws of physics, all of them in one form. The second volume asks what new results the same language can predict, beginning with the Coupling Force.

A word on honesty. This book presents a proposal. Some of its statements are definitions that we choose, some are results that follow from those choices, and some are ideas that still wait for the verdict of measurement. I have tried to label each one clearly, so that you can always tell which is which, and so that you can judge the proposal for yourself. That is what the reader of a new idea deserves.

\placeholder{A personal note from the author on how the journey began.}

\bigskip

\noindent\spacedlowsmallcaps{Acknowledgments}

\medskip

\noindent I gratefully acknowledge the support of Morocco Aerospace Industries S.A.S. \placeholder{Further acknowledgments.}

\bigskip

\begin{flushright}
\myName\\
\placeholder{Place}, \placeholder{Month} \myTime
\end{flushright}
